%=====================================================================
%  SECCION 4 -- ANALISIS DE IMPACTO Y PRIVACIDAD DESDE EL DISENO
%  Responsable: Fabian Moreno Ugarte
%
%  GENERADO desde secciones/04_analisis_impacto.md, que sigue siendo la
%  FUENTE de esta seccion. No editar aqui: editar el .md y reconvertir.
%
%  Los marcadores label reproducen los anclajes de la version archivada. La
%  numeracion 4.1-4.5 es vinculante.
%=====================================================================

\section{Análisis de impacto y privacidad desde el diseño}
\label{sec:analisis}

% ---- 4.1 ----
\subsection{Metodología del análisis de impacto}
\label{sec:eipd}

Esta sección evalúa el impacto del sistema propuesto sobre los derechos
de las personas observadas y sobre el entorno en que se despliega. Su
estructura no es arbitraria: reproduce deliberadamente la de una
Evaluación de Impacto en Protección de Datos, esto es, identificar el
tratamiento, valorar su necesidad y proporcionalidad, identificar los
riesgos para los derechos de los titulares y proponer medidas de
mitigación.

Esa correspondencia no es una analogía metodológica. Como se expuso en
la Sección 3.4, el D.S. 016-2024-JUS exige una evaluación de impacto
antes de iniciar tratamientos de alto riesgo, y la videovigilancia
masiva figura expresamente entre los supuestos que la activan. El
sistema parte de la captación masiva y continua de imágenes de personas
en las zonas comunes de un terminal aeroportuario, de modo que el
supuesto se activa con independencia de que el modelo produzca conteos
agregados y no identifique a nadie. Esta sección es, por tanto, el
insumo sustantivo de un documento que el proyecto está obligado a
producir, y no un ejercicio analítico de acompañamiento.

Conviene anticipar una tentación previsible y descartarla expresamente:
argumentar que, como la salida es agregada y no hay reconocimiento
facial, el sistema no alcanza el umbral de alto riesgo y la evaluación
no sería exigible. Este informe recomienda no construir el argumento en
esa dirección, por dos razones. La primera es que el supuesto que se
activa es el de captación masiva, y esa captación existe con
independencia de lo que el modelo haga después. La segunda es que la
evaluación de impacto es precisamente el documento en el que las
decisiones de diseño de este informe (minimización, no persistencia,
salida agregada, procesamiento local) se registran como medidas de
mitigación y quedan acreditadas ante la autoridad. Renunciar a ella
sería renunciar al lugar donde el trabajo de cumplimiento se vuelve
exigible y demostrable.

Un criterio gobierna la redacción de todo lo que sigue: se prefiere
declarar un impacto como no cuantificado antes que estimarlo sin base.
Varias magnitudes que naturalmente irían aquí (número de pasajeros
afectados, costo evitado, consumo energético, tasa de error por grupo
demográfico) dependen de información operativa que solo LAP posee o de
mediciones que el equipo aún no ha realizado. Donde ocurre, se indica
expresamente en lugar de rellenarse con una cifra plausible.

% ---- 4.2 ----
\subsection{Sesgos del sistema}

El sistema puede fallar de manera desigual, y conviene distinguir tres
fuentes de desigualdad que operan por mecanismos distintos y admiten
mitigaciones distintas. En los tres casos se indica qué reconocen los
propios autores de la literatura revisada, porque la solidez del
análisis depende de no atribuirles limitaciones que no declararon ni de
silenciar las que sí.

\textbf{Brecha de dominio.} Es la limitación mejor documentada y la más
grave para esta fase del proyecto. Los autores de P2R
\cite{Lin_2025_CVPR} reconocen expresamente que el desempeño de los
métodos de conteo se degrada de forma marcada cuando el entorno de
despliegue difiere del de entrenamiento, y que sin una etapa de
adaptación de dominio las cifras publicadas no se trasladan. Documentan
además un modo de fallo concreto del paradigma punto a punto: el
vecindario de cada píxel de primer plano queda sobre-activado y se lee
como instancia distinta, lo que produce sobreestimación severa del
conteo. La consecuencia para el proyecto es directa y condiciona la
lectura de toda la Sección de Resultados: esta primera entrega se
desarrolla sobre conjuntos públicos, sin acceso todavía a material del
Aeropuerto Jorge Chávez, de modo que ninguna cifra que se reporte debe
leerse como predicción de exactitud para el terminal. Es, en el mejor de
los casos, una cota de referencia sobre datos públicos, sujeta a
validación en sitio antes de comunicarse a LAP como expectativa de
desempeño.

\textbf{Sesgo demográfico.} Aquí la situación es distinta y debe
enunciarse con precisión: no es que los autores reconozcan un sesgo
demográfico, sino que ninguno de los trabajos revisados reporta
desempeño desagregado por características demográficas de las personas
contadas. No se trata de una omisión de este informe sino de una
ausencia en la literatura consultada, y tiene consecuencias concretas en
un aeropuerto internacional, donde la población observada es
considerablemente más heterogénea que la de los conjuntos de
entrenamiento disponibles. Los factores que razonablemente podrían
inducir desempeño desigual son identificables aunque no estén
cuantificados. La estatura es el primero: un contador entrenado sobre
anotaciones de cabezas puede subcontar sistemáticamente a niños, sobre
todo en escenas densas donde quedan ocluidos por adultos. La vestimenta
y los cubrimientos de cabeza (sombreros, gorras, capuchas, prendas
religiosas) alteran la apariencia de la región de la cabeza, que es
precisamente la señal que estos métodos explotan. El equipaje voluminoso
y las sillas de ruedas presentan siluetas distintas de las típicas en
los conjuntos de origen. Y la composición demográfica de esos conjuntos,
procedentes mayoritariamente de escenas de otras regiones del mundo, no
ofrece garantía alguna de neutralidad respecto de la población
retratada: la brecha documentada en \cite{Lin_2025_CVPR} es geográfica y
de escena, y no hay evidencia de que sea neutral respecto de quién
aparece en la imagen.

\textbf{Sesgo de anotación.} Es el menos visible de los tres y el que
los autores sí documentan de forma explícita. Los responsables de P2PNet
\cite{Song_2021_ICCV} señalan que su métrica de localización más
exigente, el nAP con umbral 0,05, cae de forma notable por efecto de las
desviaciones de etiquetado, es decir, por la imprecisión con que los
anotadores humanos sitúan el punto sobre la cabeza. El sesgo no procede
aquí del modelo sino de la verdad de referencia contra la que se lo
entrena y se lo mide. En la misma línea, los autores de CSRNet
\cite{Li_2018_CVPR} construyen su mapa de referencia difuminando cada
anotación con un núcleo gaussiano geométricamente adaptativo, lo que
incorpora un supuesto sobre la relación entre distancia a la cámara y
tamaño aparente de la cabeza; cuando ese supuesto no se cumple
(perspectivas oblicuas, techos altos, cámaras de gran angular, todas
ellas habituales en un terminal), la referencia misma queda sesgada. Los
propios autores reconocen además que sobre UCSD, un conjunto de escenas
relativamente dispersas, superan a la mayoría de métodos previos excepto
a MCNN en MAE, y que la baja resolución de esos cuadros les obligó a
reescalar por interpolación bilineal. Esto último importa operativamente
más de lo que parece: el régimen de baja densidad no es el terreno donde
el método brilla, y buena parte del día un terminal aéreo opera
precisamente en baja densidad.

Conviene ser preciso sobre la gravedad de estos sesgos, para no
exagerarla ni minimizarla. El sistema no toma decisiones sobre personas
individuales: no deniega embarques ni señala sospechosos, de modo que un
error de conteo no perjudica directamente a la persona mal contada. Sin
embargo, si el subconteo se concentra en un grupo, las zonas del
terminal donde ese grupo se concentra reciben sistemáticamente menos
recursos: menos personal y menos apertura de mostradores. El daño no es
individual sino distributivo, y es real aunque sea indirecto. De ahí se
sigue una recomendación concreta: que la campaña de validación en sitio
incluya una evaluación de desempeño desagregada, al menos por franja de
estatura y por presencia de equipaje voluminoso. La recomendación es
condicional en sentido estricto, porque sin acceso a material del
terminal no es realizable, y esa dependencia refuerza la solicitud de
acceso planteada en la Sección 3.4.

% ---- 4.3 ----
\subsection{Sesgo de automatización y factor humano}
\label{sec:impacto-social}

Un riesgo distinto de los anteriores, y menos discutido, es la
sustitución progresiva del criterio humano por el del sistema. Si el
personal de operaciones llega a confiar en el tablero más que en su
propia observación, un fallo del sistema en un momento crítico produce
un efecto peor que su ausencia, porque la atención humana ya se retiró.
El fenómeno no depende de que el modelo funcione mal: depende de que
funcione bien durante el tiempo suficiente para que nadie lo cuestione.

Este riesgo tiene un anclaje normativo directo. Como se expuso en la
Sección 3.3, el artículo 23 de la Ley 29733 reconoce el derecho al
tratamiento objetivo, esto es, a no verse sometido a una decisión con
efectos jurídicos o significativos sustentada únicamente en un
tratamiento automatizado. La medida M-10, conforme a la cual el sistema
emite alertas y no ejecuta acciones automáticas sobre personas, responde
a esa exigencia y mantiene al sistema fuera del supuesto del artículo
23. Pero M-10 es insuficiente por sí sola, porque describe la
arquitectura y no la práctica: un operador que aprueba mecánicamente
cuanto el tablero propone reintroduce de hecho la decisión automatizada
que la medida pretendía evitar, aunque formalmente medie una acción
humana.

De ahí que la capacitación del personal deba tratarse como medida de
mitigación y no como accesorio de implantación. Los operadores necesitan
conocer las limitaciones del modelo, y en particular tres de las
documentadas en esta sección: que el desempeño sobre datos públicos no
predice el desempeño en el terminal mientras no exista validación en
sitio; que existe un modo de fallo por sobreestimación severa
documentado en \cite{Lin_2025_CVPR}; y que el régimen de baja densidad,
habitual durante buena parte de la jornada, no es aquel en que estos
métodos se comportan mejor. Un operador que conoce estos tres puntos
interpreta una alerta anómala como lo que probablemente es; uno que no
los conoce, la ejecuta.

Existe además un impacto sobre el comportamiento de las personas
observadas que conviene dejar planteado con la cautela que corresponde.
Se sostiene habitualmente que saberse observado modifica la conducta;
este informe no puede afirmarlo como hecho acreditado, porque las
referencias validadas del proyecto son los artículos técnicos de
antecedentes y no se ha incorporado literatura sobre efectos
conductuales de la videovigilancia. Formulado en condicional: si ese
efecto existe en la magnitud que la discusión pública le atribuye, en un
aeropuerto, donde la vigilancia ya está normalizada, el efecto marginal
de añadir un sistema de conteo sería previsiblemente pequeño, aunque no
nulo, y la normalización misma formaría parte del fenómeno. Sostener el
punto con firmeza exigiría respaldo bibliográfico que este informe
todavía no tiene.

% ---- 4.4 ----
\subsection{Lectura ambiental del costo de cómputo}
\label{sec:impacto-ambiental}

El entrenamiento de modelos de visión es intensivo en cómputo y, por
tanto, en energía. El equipo no ha medido el consumo energético del
entrenamiento ni el de la inferencia, de modo que lo que sigue es
cualitativo y estructural; se indica en cada caso qué medición
permitiría convertirlo en cuantitativo.

Tres decisiones de diseño ya adoptadas o previstas tienen efecto directo
sobre la huella del entrenamiento. La primera es la reutilización de
modelos preentrenados: CSRNet parte de las primeras capas
convolucionales de VGG-16 ya entrenadas \cite{Li_2018_CVPR}, lo que
evita entrenar desde cero. La segunda es la eficiencia de etiquetado del
enfoque semi-supervisado de \cite{Lin_2025_CVPR}, sobre la que conviene
no sobrevender el argumento: reduce el esfuerzo humano de anotación, no
necesariamente el cómputo, puesto que el esquema maestro-estudiante
mantiene dos modelos en entrenamiento.

La tercera es la que admite una lectura ambiental más precisa, porque
los propios autores la cuantifican en prosa. P2R prescinde del algoritmo
húngaro, componente que los autores identifican expresamente como
costoso en tiempo dentro de la estrategia de emparejamiento punto a
punto. Sobre una imagen de 576 × 960 píxeles con 775 puntos anotados, el
cálculo de la pérdida punto a punto requiere 0,4307 segundos frente a
los 0,0064 segundos de P2R, cerca de sesenta y ocho veces más rápido
\cite{Lin_2025_CVPR}. La lectura ambiental de ese dato conviene hacerla
con cuidado, porque se presta a exageración: la cifra mide el costo de
una operación dentro de cada iteración de entrenamiento, no el consumo
total del ciclo, y el ahorro agregado depende del número de iteraciones
y de qué fracción del tiempo por iteración representaba esa operación.
Lo que sí puede afirmarse sin exceder el dato es que la elección de la
función de pérdida no es ambientalmente neutra, y que en un
entrenamiento con adaptación de dominio en sitio, que es el que la
brecha documentada obliga a contemplar, una diferencia de ese orden en
un componente ejecutado en cada iteración es un factor a considerar
junto con la exactitud.

A escala de la vida útil del sistema, el consumo de la inferencia
continua supera previsiblemente al del entrenamiento: este ocurre unas
pocas veces, aquella se ejecuta permanentemente sobre múltiples flujos
de video. Dos decisiones de diseño con motivación distinta convergen en
reducirlo. La frecuencia de muestreo es la primera: las aglomeraciones
se forman en escalas de minutos, de modo que procesar un fotograma cada
varios segundos en lugar de a velocidad completa reduce el cómputo en un
factor elevado sin pérdida operativa relevante, y reduce simultáneamente
el volumen de imágenes tratadas. La segunda es el procesamiento en el
borde, adoptado como medida M-01 por razones de protección de datos, que
evita transmitir video de forma continua hacia la nube con el ahorro de
transmisión asociado.

Este es el resultado analítico más útil de la subsección y conviene
presentarlo como tal ante el cliente: las decisiones que reducen la
exposición de datos personales (muestrear menos, procesar localmente, no
persistir imágenes, agregar la salida) son exactamente las mismas que
reducen el consumo de cómputo, de transmisión y de almacenamiento. Ambos
principios apuntan en la misma dirección, lo que refuerza el argumento
de privacidad con una justificación adicional que no depende del
cumplimiento normativo.

Quedan dos puntos abiertos. El primero es que convertir esta subsección
de cualitativa en cuantitativa depende de una sola medición: las horas
de GPU efectivamente empleadas durante el desarrollo. Es un registro
sencillo, pero no es reconstruible a posteriori: si no se lleva desde el
inicio del entrenamiento, la huella del desarrollo queda definitivamente
sin cuantificar. Por esa razón, este informe recomienda registrar las
horas de GPU desde el primer entrenamiento. El segundo es la huella de
hardware: si el proyecto requiere cámaras o servidores nuevos, su
fabricación y su disposición final integran el balance, y reutilizar la
infraestructura de CCTV existente es, además de la opción económicamente
decisiva, la ambientalmente preferible. Se recomienda, por último,
verificar si LAP cuenta con una política institucional de sostenibilidad
o con compromisos de reporte ambiental, porque, de existir, el análisis
debería alinearse con sus categorías de medición en lugar de proponer un
marco paralelo.

% ---- 4.5 ----
\subsection{Riesgos residuales y decisiones pendientes}
\label{sec:sintesis-impacto}
\label{sec:impacto-etico}

Integrando lo anterior, el perfil de riesgo del sistema admite una
síntesis en cuatro dimensiones. En la ética, el impacto positivo es
ofrecer una alternativa menos invasiva que la identificación biométrica
para la misma finalidad, y el riesgo principal es el deslizamiento hacia
identificación o seguimiento, mitigable mediante la declaración de
finalidades excluidas y las medidas M-02 y M-03. En la social, el
beneficio es la prevención de incidentes por aglomeración y la
asignación de recursos basada en evidencia, y los riesgos son la
distribución desigual de beneficio y carga, el uso laboral no declarado
y el exceso de confianza en el sistema. En la económica, el beneficio es
la eficiencia operativa y el valor de los incidentes evitados, y el
riesgo dominante es el costo de reconstruir el sistema si se detecta una
no conformidad después del despliegue, lo que hace de la privacidad
desde el diseño la opción económicamente racional además de la
obligatoria. En la ambiental, el beneficio es una huella menor que la de
alternativas de mayor granularidad, y el riesgo es el consumo acumulado
de la inferencia continua a lo largo de la vida útil.

Tres conclusiones transversales se desprenden del análisis. La primera
es que las decisiones de minimización operan simultáneamente en las
cuatro dimensiones: no persistir imágenes, agregar la salida y muestrear
a baja frecuencia reducen a la vez el riesgo ético, la exposición
social, el costo de almacenamiento y la huella ambiental. Son el núcleo
de la propuesta de este informe precisamente porque ningún objetivo
compite con otro.

La segunda es que el riesgo dominante es de gobernanza y no de
tecnología. El deslizamiento de finalidad y el uso no declarado no se
previenen con mejor ingeniería, sino con documentación, compromisos
contractuales y auditoría. Esto significa que la documentación de
cumplimiento no es accesoria al proyecto técnico: es la mitigación
principal del riesgo principal, y así debería exponerse ante el cliente.

La tercera es que varios impactos permanecen sin cuantificar por falta
de información del cliente y de mediciones propias, y se ha optado por
declararlo antes que por estimar sin base.

El riesgo de deslizamiento de finalidad merece un desarrollo propio,
porque es el principal y no proviene de un defecto técnico sino del
éxito del sistema. Una vez desplegada la infraestructura (cámaras
calibradas, canal de video, capacidad de inferencia, tablero de
visualización), el costo marginal de añadir capacidades es bajo:
incorporar reidentificación para medir tiempos de permanencia, o
reconocimiento facial para detectar personas de interés, es una decisión
de configuración y de presupuesto, no un rediseño. La barrera que impide
ese deslizamiento no es técnica sino documental y contractual. La
mitigación que este informe propone es documentar la finalidad de forma
restrictiva e incorporar al entregable una declaración de finalidades
excluidas: el sistema no realiza identificación de personas, no realiza
seguimiento individual de trayectorias, no realiza reconocimiento de
emociones y no alimenta decisiones automatizadas sobre pasajeros
individuales. Una lista de lo que el sistema \emph{no} hace es más
difícil de erosionar silenciosamente que una lista de lo que sí hace, y
su valor depende de que sea explícita: una finalidad excluida que solo
consta en el ánimo del equipo no restringe nada. A esa declaración
debería incorporarse expresamente la exclusión del uso laboral de los
datos, porque si los conteos por zona se cruzan con los turnos del
personal, el sistema se convierte de facto en una herramienta de
evaluación de desempeño, con implicancias laborales propias y con una
finalidad distinta de la declarada.

Los riesgos residuales que este informe no cierra son de tres clases, y
conviene no presentarlos como si fueran del mismo tipo. Los primeros son
ausencias de evidencia subsanables dentro del proyecto si se obtiene
acceso a material del terminal: la validación de exactitud en el dominio
de destino y la evaluación de sesgo demográfico. Los segundos son
decisiones del equipo que siguen abiertas y de las que depende buena
parte de este análisis: la arquitectura definitiva y la granularidad de
salida, la localización del procesamiento (\emph{on-premise} o nube
externa), el plazo concreto de retención y el registro de horas de GPU.
Los terceros son discrepancias entre documentos del propio proyecto, que
este informe deja registradas y expresamente no resuelve por su cuenta.
Son tres, y no todas tienen el mismo signo. El uso del término «personas
anonimizadas» para describir lo que la Ley 29733 califica como
disociación aparece tanto en la presentación preliminar del proyecto
como en la documentación del módulo de almacenamiento del prototipo, y
está tratado en la Sección 3.2. La medición de permanencia en zona
comercial como finalidad distinta de la estimación de aglomeraciones,
que exigiría base de legitimación propia, se trata en la Sección 3.2; el
esquema de datos del prototipo, que registra zona y marcas de primera y
última detección por persona, la hace técnicamente posible hoy. En
cambio, la preocupación sobre reidentificación mediante
\emph{embeddings} de apariencia \textbf{no se confirmó}: el prototipo
resuelve la fusión entre cámaras por asociación espacio-temporal y
reserva un descriptor de color no biométrico como mero desempate, de
modo que la decisión de diseño va en la dirección que exige el requisito
de no derivar de la imagen rasgos que permitan identificar a una
persona, ni siquiera de forma indirecta, según se detalla en la Sección
3.4.

A esas discrepancias se añade un riesgo residual de distinta clase,
porque no es una divergencia documental sino un incumplimiento
verificable: \textbf{la medida M-03 no se cumple en la implementación
actual}, que persiste coordenadas individuales enlazadas a un
identificador por persona en lugar de agregar por zona. Es el único
hallazgo de todo este análisis que el equipo puede cerrar por sí solo,
de inmediato y sin depender de LAP ni de acceso a datos del terminal, y
por eso debería encabezar la lista de acciones. Su cierre depende
únicamente del equipo, que puede agregar por zona antes de escribir o
fijar una retención muy corta sobre la tabla de posiciones; mientras no
lo haga, el cumplimiento de M-03 no puede afirmarse ante el cliente.
